\documentclass[11pt]{article}
\usepackage{amsmath,amsthm,amssymb}
\usepackage[margin=1in]{geometry}
\usepackage{hyperref}
\usepackage{booktabs}
\usepackage{enumitem}

\newtheorem{theorem}{Theorem}
\newtheorem{lemma}[theorem]{Lemma}
\newtheorem{proposition}[theorem]{Proposition}
\newtheorem{corollary}[theorem]{Corollary}
\theoremstyle{definition}
\newtheorem{definition}[theorem]{Definition}
\newtheorem{remark}[theorem]{Remark}
\newtheorem{observation}[theorem]{Observation}

\newcommand{\Hq}{H_q}
\newcommand{\Sn}{S_n}
\newcommand{\Vlam}{V_\lambda}
\newcommand{\Bdec}{B_{\ell+1}^{(\lambda)}}
\newcommand{\Om}{\Omega}

\title{Structural skeleton of $\Bdec \Vlam$:\\
  the leftmost-and-rightmost-factor theorem}
\author{Clio}
\date{13 May 2026 (evening)}

\begin{document}
\maketitle

\begin{abstract}
The May-13 morning paper refuted the naive Hecke-iso upgrade of the
dim-formula
$\dim \Bdec V_\lambda = f^{\lambda^{(\ge 2)}}$: no parabolic embedding
$H_q(S_{n-\ell}) \hookrightarrow H_q(\Sn)$ acts on $B$. Computationally,
the only generator that \emph{structurally} acts as a scalar on $B$ was
$T_\ell$ (acting as $q$), with shape-dependent additions
$T_1, \ldots, T_{k(\lambda)}$ acting as $-1$.

We prove the two universal facts here, both with one-line proofs from
the Hecke relations:

\medskip
\textbf{Theorem A.} $T_\ell \cdot \Om = q \cdot \Om$ in $\Hq$, where
$\Om := R'_{\ell+1} R'_{\ell+2} \cdots R'_n$. Consequently $T_\ell$ acts
as scalar $q$ on $\Bdec V_\lambda$, i.e.\ $\Bdec V_\lambda \subseteq
E_\ell^+|_{V_\lambda}$.

\medskip
\textbf{Theorem A$'$.} $\Om \cdot T_1 = q \cdot \Om$ in $\Hq$.
Consequently $\Om$ annihilates the $(-1)$-eigenspace
$E_1^- = \ker(T_1+1)|_{V_\lambda}$, and the induced map
$V_\lambda / E_1^- \to E_\ell^+$ has image $\Bdec V_\lambda$.

\medskip
The two theorems express the same Hecke identity
$(T_j-q)(T_j+1)=0$ applied at the \emph{leftmost} and \emph{rightmost}
factor of the chain product $\Om$. They are dual structural skeletons
of $\Om$ and, together, the only \emph{universal} eigenspace
constraints on $\Bdec V_\lambda$: everything else (e.g.\ which
$T_j$, $j < \ell$, acts as $-1$) is shape-dependent.

The result tightens but does not contradict the May-13 morning refutation:
$T_\ell$ acting as scalar~$q$ is the strongest \emph{shape-uniform}
Hecke-equivariance statement available, and we explain precisely why no
analogous statement holds for $T_j$ with $j \notin \{\ell\}$ or
$j \notin \{1, \ldots, k(\lambda)\}$.

Verified mod $P = 100\,003$ at $q = 11$ for $\lambda \in \{(3,2,1^3),
(4,2,1^4), (3,3,1^4), (4,3,1^3), (3,2,1^4), (2,1^3), (3,1^3), (3,2)\}$.
\end{abstract}

\section{Setup}\label{sec:setup}

We retain the notation of the May-13 multi-corner dimension paper.
$\Hq = \Hq(\Sn)$ is the Iwahori--Hecke algebra with generators
$T_1, \ldots, T_{n-1}$ and relations
\begin{equation}\label{eq:hecke-rel}
   T_j^2 = (q-1) T_j + q, \qquad
   T_j T_{j+1} T_j = T_{j+1} T_j T_{j+1}, \qquad
   T_i T_j = T_j T_i \;\;(|i-j| \ge 2).
\end{equation}
For $j \ge 1$, set
\begin{equation}\label{eq:Sj}
   S_j \;:=\; T_j + 1 \in \Hq.
\end{equation}
The Hecke relation
\eqref{eq:hecke-rel} factors as $(T_j - q)(T_j+1) = 0$. Equivalently
\begin{equation}\label{eq:absorption}
   T_j S_j \;=\; q\, S_j, \qquad S_j T_j \;=\; q\, S_j.
\end{equation}
That is: $S_j$ \emph{absorbs} $T_j$ from either side, producing
scalar $q$.

Let $V_\lambda$ be the Specht module in the Hoefsmit seminormal basis,
indexed by SYTs of $\lambda$. Eigenspaces of $T_j$ in $V_\lambda$:
\[
   E_j^+ \;:=\; \ker(T_j - q)|_{V_\lambda}, \qquad
   E_j^- \;:=\; \ker(T_j + 1)|_{V_\lambda}.
\]
Note $E_j^+ = \mathrm{Im}(S_j|_{V_\lambda})$ and
$E_j^- = \ker(S_j|_{V_\lambda})$.

For $k \ge 1$, define the level-$k$ descent operator
\[
   R'_k \;:=\; S_{k-1} S_{k-2} \cdots S_2 S_1 \;=\;
   (T_{k-1}+1)(T_{k-2}+1)\cdots(T_2+1)(T_1+1) \;\in\; \Hq.
\]
For $\lambda \vdash n$ and $0 \le \ell \le n-1$, set
\begin{equation}\label{eq:Omega}
   \Om \;=\; \Om^{(\lambda)}_\ell \;:=\;
   R'_{\ell+1} R'_{\ell+2} \cdots R'_n \;\in\; \Hq,
\end{equation}
the chain product from level $\ell+1$ down to level $n$. We are
interested in the iterated image
\[
   B \;=\; \Bdec V_\lambda \;:=\; \Om \cdot V_\lambda \;\subseteq\; V_\lambda.
\]
Throughout, $\ell$ is the length of $\lambda$ (so $\Om$ takes the
``decoration-shape'' product), but the structural theorems do not require
this. Set $\Bdec V_\lambda := V_\lambda$ if $\ell+1 > n$ (empty product).

\section{The two structural skeleton theorems}\label{sec:two-theorems}

The Hecke absorption identity \eqref{eq:absorption} drives both results.
Read \eqref{eq:Omega} from left to right: the \emph{leftmost} factor of
$\Om$ is $S_\ell$ (the leading $S$ in $R'_{\ell+1}$). The \emph{rightmost}
factor of $\Om$ is $S_1$ (the trailing $S$ in $R'_n$).

\begin{theorem}[A: $T_\ell$ acts as scalar $q$]\label{thm:A}
For any partition $\lambda \vdash n$ and any $\ell$ with $1 \le \ell \le n-1$,
\[
   T_\ell \cdot \Om \;=\; q \cdot \Om \qquad \text{in } \Hq.
\]
In particular, $T_\ell$ acts as scalar $q$ on $\Bdec V_\lambda$,
i.e.\ $\Bdec V_\lambda \subseteq E_\ell^+|_{V_\lambda}$.
\end{theorem}

\begin{proof}
The leftmost factor of $R'_{\ell+1}$ is $S_\ell$:
$R'_{\ell+1} = S_\ell \cdot (S_{\ell-1} \cdots S_1) = S_\ell \cdot R'_\ell$.
Hence
\[
   \Om \;=\; R'_{\ell+1} \cdot (R'_{\ell+2} \cdots R'_n)
       \;=\; S_\ell \cdot R'_\ell \cdot (R'_{\ell+2} \cdots R'_n)
       \;=\; S_\ell \cdot X
\]
for $X := R'_\ell \cdot R'_{\ell+2} \cdots R'_n \in \Hq$. By
\eqref{eq:absorption},
$T_\ell \Om = T_\ell S_\ell X = q S_\ell X = q \Om$. Applying to any
$v \in V_\lambda$ and writing $b = \Om v \in B$:
$T_\ell b = T_\ell \Om v = q \Om v = q b$. \qed
\end{proof}

\begin{theorem}[A$'$: $\Om$ annihilates $E_1^-$]\label{thm:A-prime}
For any partition $\lambda \vdash n$ and any $\ell$ with $0 \le \ell \le n-1$,
\[
   \Om \cdot T_1 \;=\; q \cdot \Om \qquad \text{in } \Hq.
\]
Equivalently, $\Om$ annihilates $E_1^- = \ker(T_1+1)|_{V_\lambda}$, and
the induced map $\bar\Om: V_\lambda / E_1^- \to V_\lambda$ has image
$\Bdec V_\lambda$.
\end{theorem}

\begin{proof}
The rightmost factor of $R'_n$ is $S_1$:
$R'_n = (S_{n-1} \cdots S_2) \cdot S_1$. Hence
\[
   \Om \;=\; (R'_{\ell+1} \cdots R'_{n-1}) \cdot R'_n
       \;=\; (R'_{\ell+1} \cdots R'_{n-1}) \cdot (S_{n-1} \cdots S_2) \cdot S_1
       \;=\; Y \cdot S_1
\]
for $Y \in \Hq$. By \eqref{eq:absorption},
$\Om T_1 = Y S_1 T_1 = q Y S_1 = q \Om$.

For the equivalence: for any $v \in V_\lambda$, write $\Om v$. If
$v \in E_1^-$, then $T_1 v = -v$, so $S_1 v = (T_1+1) v = 0$, hence
$\Om v = Y S_1 v = 0$.

Conversely, the identity $\Om(T_1 v) = q \Om(v)$ holds for all $v$;
if $v \in E_1^-$ then $\Om(T_1 v) = \Om(-v) = -\Om(v)$, so
$(-1)\Om(v) = q \Om(v)$, forcing $(q+1)\Om(v) = 0$. At generic $q$
(in particular at $q = 11$ mod $P = 100\,003$, $q+1 = 12 \ne 0$),
this gives $\Om(v) = 0$. Both directions show $E_1^- \subseteq \ker(\Om|_{V_\lambda})$.
\qed
\end{proof}

\begin{corollary}[The structural skeleton]\label{cor:skeleton}
The operator $\Om|_{V_\lambda}: V_\lambda \to V_\lambda$ induces a
linear map
\[
   \bar\Om: V_\lambda / E_1^- \;\longrightarrow\; E_\ell^+
\]
whose image is $\Bdec V_\lambda$. In symbols,
\[
   \Bdec V_\lambda \;\subseteq\; E_\ell^+, \qquad
   \ker(\Om|_{V_\lambda}) \;\supseteq\; E_1^-.
\]
\end{corollary}

\begin{proof}
Theorem~\ref{thm:A} gives the codomain $E_\ell^+$;
Theorem~\ref{thm:A-prime} gives the kernel containment. \qed
\end{proof}

\begin{remark}[``Skeleton'' interpretation]
The two theorems are dual structural facts: the leftmost factor of $\Om$
forces $T_\ell$-eigenspace constraints \emph{on the image} $B$; the
rightmost factor forces $T_1$-eigenspace constraints \emph{on the
kernel}. They sit at the two ``ends'' of the chain product
$\Om = R'_{\ell+1} R'_{\ell+2} \cdots R'_n$ and are the only
constraints from $\Hq$ alone (without using the specific
representation $V_\lambda$).
\end{remark}

\section{What is NOT universal: shape-dependence}\label{sec:shape}

The May-13 morning refutation paper found computationally that for
$\lambda = (3,2,1^3)$ at $n = 8$, $\ell = 5$, $\dim B = 2$:
\begin{itemize}[topsep=0.2em,itemsep=0.1em]
\item $T_5$ acts as scalar $q$ on $B$ \quad (Theorem~\ref{thm:A});
\item $T_1$ acts as scalar $-1$ on $B$ \quad (shape-dependent);
\item $T_j$ for $j \in \{2,3,4,6,7\}$ does \emph{not} even preserve $B$.
\end{itemize}
For $\lambda = (4, 3, 1^3)$, the same refutation paper showed $T_1$ does
\emph{not} act as $-1$: there is no length-$k$ prefix of $\{T_1, \ldots,
T_{\ell-1}\}$ acting as scalars beyond $T_\ell$ itself.

Theorem~\ref{thm:A} explains the universal $T_\ell = q$ behaviour
structurally. What about the empirical $T_1, T_2, \ldots, T_k$ acting as
$-1$ on $B$, for some shape-dependent $k = k(\lambda) \in \{0, 1, 2\}$?

\begin{observation}[Why $T_1 = -1$ is \emph{not} universal]
The natural attempt to prove $T_1 \cdot \Om \subseteq -1 \cdot \Om$
fails because $T_1$ does \emph{not} satisfy a left-absorption identity
with $\Om$: $T_1$ commutes with $T_j$ for $j \ge 3$ but not with $T_2$,
and $T_2$ appears (as $S_2$) in every $R'_k$ for $k \ge 3$. So
$T_1$ cannot be pulled across $\Om$ to a position where the Hecke
relation gives a scalar.

What \emph{is} structurally true is the \emph{right}-action identity:
$\Om T_1 = q \Om$ (Theorem~\ref{thm:A-prime}). This says $\Om$
acts as zero on $E_1^-$ (in the \emph{source}), not that $B$ lies
in $E_1^-$ (in the \emph{target}). The image $\Om(E_1^+) = B$ may or
may not lie in $E_1^-$, depending on whether the operator
$\Om|_{E_1^+}: E_1^+ \to V_\lambda$ happens to factor through $E_1^-$.
That is a shape-dependent condition. Empirically it holds for
$\lambda = (3,2,1^3)$ and many other shapes; for $\lambda = (4,3,1^3)$
it does not.
\end{observation}

The structural skeleton (Corollary~\ref{cor:skeleton}) is the strongest
\emph{shape-uniform} statement available. Everything beyond it is
shape-dependent, and that is precisely what the May-13 morning
refutation paper made explicit.

\section{Generalisation to all levels of the chain}\label{sec:chain-levels}

The same proof of Theorem~\ref{thm:A} applies at every level of the
chain product, by replacing $\ell$ with any $j$ in the construction.
Define
\[
   B_{j+1} \;=\; B_{j+1}^{(\lambda)} V_\lambda \;:=\;
   R'_{j+1} R'_{j+2} \cdots R'_n \cdot V_\lambda, \qquad 0 \le j \le n-1.
\]
(So $B_{j+1}$ is the chain product starting from level $j+1$.)

\begin{theorem}[Level-$j$ scalar fact]\label{thm:level-j}
For every $j \in \{1, 2, \ldots, n-1\}$ and every $\lambda \vdash n$:
\[
   T_j \cdot \big(R'_{j+1} R'_{j+2} \cdots R'_n\big) \;=\; q \cdot
   \big(R'_{j+1} R'_{j+2} \cdots R'_n\big) \qquad \text{in } \Hq.
\]
Hence $T_j$ acts as scalar $q$ on $B_{j+1}^{(\lambda)} V_\lambda$, i.e.\
$B_{j+1}^{(\lambda)} V_\lambda \subseteq E_j^+|_{V_\lambda}$.
\end{theorem}

\begin{proof}
Same as Theorem~\ref{thm:A}, with $\ell$ replaced by $j$. \qed
\end{proof}

\begin{corollary}[Eigenspace chain]\label{cor:eigenspace-chain}
The descending chain $B_n \supseteq B_{n-1} \supseteq \cdots \supseteq B_{\ell+1}$
need not be literally nested (each $B_{j+1}$ is the image of $R'_{j+1}$
on the next), but each $B_{j+1}$ is contained in $E_j^+$:
\[
   B_{j+1} \;\subseteq\; E_j^+ \quad\text{for all } 1 \le j \le n-1.
\]
\end{corollary}

\begin{remark}
This sharpens the May-13 morning Step-1 conjecture (\textsc{Prove.md},
``$B$ is fixed under $(T_{j-1}+1)$ for $j > \ell$''). The correct
universal statement is: $B_{j+1}$ is in the $+q$-eigenspace of $T_j$,
not in any joint eigenspace involving $T_{j+1}, T_{j+2}, \ldots$.
The latter would have been the Hecke-iso statement that was refuted.
\end{remark}

\section{Computational verification}\label{sec:verification}

We verify Theorem~\ref{thm:A} (specifically, $T_\ell \cdot W = q \cdot W$
where $W$ is an explicit column basis of $B$) and Theorem~\ref{thm:A-prime}
(specifically, $\Om \cdot E_1^- = 0$ as a matrix equation) mod
$P = 100\,003$ at $q = 11$, using the Hoefsmit seminormal matrix model
(script \texttt{compute\_dim\_B.py} of May-13).

\begin{center}
\renewcommand{\arraystretch}{1.15}
\begin{tabular}{l|c|c|c|c|c}
\toprule
$\lambda$ & $n$ & $\ell$ & $\dim V_\lambda$ & $\dim B$ & $T_\ell = q$ on $B$?\quad $\Om(E_1^-) = 0$? \\
\midrule
$(3,2)$        & 5  & 2 & 5   & 2 & yes \quad yes \\
$(2,1,1,1)$    & 5  & 4 & 4   & 1 & yes \quad yes \\
$(3,1,1,1)$    & 6  & 4 & 10  & 1 & yes \quad yes \\
$(3,2,1^3)$    & 8  & 5 & 64  & 2 & yes \quad yes \\
$(3,2,1^4)$    & 9  & 6 & 105 & 2 & yes \quad yes \\
$(4,2,1^4)$    & 10 & 6 & 350 & 3 & yes \quad yes \\
$(3,3,1^4)$    & 10 & 6 & 225 & 2 & yes \quad yes \\
$(4,3,1^3)$    & 10 & 5 & 525 & 5 & yes \quad yes \\
\bottomrule
\end{tabular}
\end{center}

\noindent (Scripts:
\texttt{2026-05-13-evening/verify\_T\_ell\_scalar.py},
\texttt{2026-05-13-evening/verify\_omega\_kills\_E1minus.py}.)

In all $8$ shapes, both structural facts hold, and the $\dim B$
column matches the May-13 morning closed form $f^{\lambda^{(\ge 2)}}$
where applicable. The proofs of Theorems~\ref{thm:A} and
\ref{thm:A-prime} are themselves Hecke-algebra identities, so they
hold over $\mathbb{Z}[q]$ uniformly; the computational table is a
sanity check, not part of the proof.

\section{What this gives the iso programme}\label{sec:iso-prog}

The May-13 morning refutation paper identified three resolutions for the
dim-equality $\dim B = f^{\lambda^{(\ge 2)}}$:
\begin{itemize}[topsep=0.2em,itemsep=0.1em]
\item[(a)] iso via exotic operators in $\Hq$;
\item[(b)] vector-space-only iso (combinatorial bijection);
\item[(c)] subquotient or induced Hecke structure.
\end{itemize}

The structural skeleton (Cor.~\ref{cor:skeleton}) constrains resolution
(a): any exotic operator $X \in \Hq$ that intertwines an
$H_q(S_{n-\ell})$-action on $B$ with $V_{\lambda^{(\ge 2)}}$ must
commute with $T_\ell$ (so as to preserve the $T_\ell = q$ eigenspace
property of $B$). Hence $X$ lies in the centraliser
$Z_{\Hq}(T_\ell)$, which is generated by $\{T_j : j \ne \ell-1, \ell+1\}
\cup \{T_\ell\}$ as a Hecke-style subalgebra of type
$A_{\ell-2} \times A_{n-\ell-2}$, with $T_\ell$ acting as scalar $q$.
However the refutation paper showed (computationally) that
$T_1, \ldots, T_{\ell-2}$ and $T_{\ell+2}, \ldots, T_{n-1}$ generally do
\emph{not} preserve $B$ as a subspace. So even within $Z_{\Hq}(T_\ell)$,
the generic generators fail. Resolution (a) requires \emph{non-trivial}
elements of $Z_{\Hq}(T_\ell)$, beyond the simple generators.

A natural candidate is the family of Jucys--Murphy elements
\[
   L_k \;=\; T_{k-1} \cdots T_2 T_1 \cdot T_1 T_2 \cdots T_{k-1}, \quad 1 \le k \le n,
\]
which generate a maximal commutative subalgebra of $\Hq$. The $L_k$
commute pairwise; $L_k$ commutes with $T_j$ when $j \notin \{k-1, k\}$.
In particular $L_k$ commutes with $T_\ell$ when $k \notin
\{\ell, \ell+1\}$. We have not investigated whether such $L_k$
preserve $B$ as a subspace, and leave this as an open direction for
the iso programme.

The dual structural fact (Theorem~\ref{thm:A-prime}: $\Om(E_1^-) = 0$)
gives the parallel constraint on the \emph{source}: any natural
$\Hq$-equivariant lift $\widetilde B \to B$ (e.g.\ from an induced
representation) must factor through the quotient $V_\lambda / E_1^-$.

\section{Status and open directions}\label{sec:status}

\paragraph*{Proved (this paper).}
Theorems~\ref{thm:A}, \ref{thm:A-prime}, \ref{thm:level-j}, with
one-line proofs from the Hecke identity $(T_j-q)(T_j+1)=0$. These
constitute the universal structural skeleton of
$\Bdec V_\lambda$.

\paragraph*{Already known (May-13 morning).}
\begin{itemize}[topsep=0.2em,itemsep=0.1em]
\item Dim formula $\dim B = f^{\lambda^{(\ge 2)}}$ in the fully-stable
regime (multi-corner dimension paper).
\item Refutation of naive parabolic Hecke iso (decoration-iso-failed paper).
\item Empirical $T_1, \ldots, T_{k(\lambda)}$ act as $-1$ on $B$, where
$k(\lambda) \in \{0, 1, 2\}$ is shape-dependent.
\end{itemize}

\paragraph*{Open.}
\begin{itemize}[topsep=0.2em,itemsep=0.1em]
\item Characterise the shape-dependence of $k(\lambda)$ (when does $T_j$
act as $-1$ on $B$, for $j < \ell$?).
\item Identify the precise stability threshold $q_0(\widehat\lambda)$
for $\dim B = f^{\lambda^{(\ge 2)}}$ (non-monotone in $q$ in general).
\item Investigate Jucys--Murphy elements $L_k$ with
$k \notin \{\ell, \ell+1\}$: do any commute with the projector to $B$?
This is the natural setting for resolution (a) of the iso programme.
\item Express $B$ as a subquotient / induced module (resolution (c)).
\end{itemize}

\paragraph*{Honest assessment.}
Theorems~\ref{thm:A} and \ref{thm:A-prime} are not deep. They are
one-line consequences of the basic Hecke identity. Their value is in
\emph{naming} the structural skeleton precisely and explaining why no
analogous statement holds for other $T_j$. Combined with the May-13
morning refutation, the rep-theoretic story of $\Bdec V_\lambda$ now has
a clean ``universal core + shape-dependent decoration'' structure:
\begin{itemize}[topsep=0.2em,itemsep=0.1em]
\item Universal core: $B \subseteq E_\ell^+$, $\ker(\Om) \supseteq E_1^-$.
\item Shape-dependent decoration: which of $T_1, \ldots, T_{\ell-1}$
acts as $-1$ on $B$; the stability threshold $q_0(\widehat\lambda)$;
the precise dimension when not in the stable regime.
\end{itemize}

\end{document}
